\documentclass[11pt,a4paper]{article}

% ---------- Encoding & Sprache ----------
\usepackage[utf8]{inputenc}
\usepackage[T1]{fontenc}
\usepackage[ngerman]{babel} % aktiviert automatische Silbentrennung

% ---------- Formalia gemaess Kapitel 4.3 der Aufgabenstellung ----------
% Seitenraender: oben/unten 2cm, links 2cm, rechts 2cm
\usepackage[top=2cm, bottom=2cm, left=2cm, right=2cm]{geometry}

% Schrifttyp: Arial ist proprietaer und in pdfLaTeX nicht direkt einbettbar.
% Helvetica (helvet) ist metrisch kompatibel und die uebliche pdfLaTeX-Alternative.
\usepackage[scaled]{helvet}
\renewcommand{\familydefault}{\sfdefault}

% Zeilenabstand 1,5
\usepackage{setspace}
\onehalfspacing

\usepackage{enumitem}
\setlength{\parindent}{0pt}
\setlength{\parskip}{6pt} % Absatzabstand 6pt nach Zeilenumbruch

% Fussnoten: Arial 10pt, Blocksatz
\usepackage[flushmargin]{footmisc}
\renewcommand{\footnotesize}{\fontsize{10}{12}\selectfont}

\begin{document}

% Blocksatz ist LaTeX-Standard (\justify), Silbentrennung via babel aktiv.

\begin{center}
{\fontsize{12}{14}\selectfont\textbf{Abstract}}\\[4pt]
{\fontsize{12}{14}\selectfont Datenbankentwurf für eine Buchtausch-App}\\[2pt]
\normalsize Tizian Senger -- Matrikelnummer: IU14143428 -- DLBDSPBDM01\_D
\end{center}

\vspace{6pt}

% =========================================================
{\fontsize{12}{14}\selectfont\textbf{Inhaltliche und konzeptionelle Beschreibung}}
% =========================================================

% TODO (ca. 0,75-1 Seite): Problem, Zielsetzung, Konzept und erreichtes Ergebnis
% der Buchtausch-App-Datenbank zusammenfassend darstellen.

% =========================================================
{\fontsize{12}{14}\selectfont\textbf{Making of -- technischer Breakdown}}
% =========================================================

% TODO (ca. 0,5-0,75 Seite): Nuechterner, informativer Ueberblick ueber die technische
% Herangehensweise: gewaehltes DBMS, Vorgehen ueber die drei Phasen, eingesetzte Tools.

\end{document}
